\import{algebra/ring/ideal.tex}
\import{algebra/group/quotient.tex}

\section{Congruence modulo an ideal}

Congruence modulo a two-sided ideal is literally the coset relation in the
additive group.  Its chosen orientation says that $x$ is congruent to $y$
when $-x+y$ lies in the ideal; additive commutativity gives the customary
equivalent condition $x-y\in I$.

\begin{definition}\label{ideal_congruence}
    $\idealCongruence{R}{I}=
        \subgroupCosetRelation{\AdditiveGroup{R}}{I}$.
\end{definition}

\begin{theorem}\label{ideal_congruence_equivalence}
    Suppose $I$ is a two-sided ideal of $R$.
    Then $\idealCongruence{R}{I}$ is an equivalence on $\carrier[R]$.
\end{theorem}
\begin{proof}
    Let $G=\AdditiveGroup{R}$.
    $\subgroup{I}{G}$ by
        \obv{two_sided_ideal,left_ideal_additive_subgroup}.
    Hence $\subgroupCosetRelation{G}{I}$ is an equivalence on
        $\carrier[G]$ by \obv{subgroup_coset_relation_equivalence}.
    $\carrier[G]=\carrier[R]$ by \auto{additive_group_fields}.
    Follows by \auto{ideal_congruence}.
\end{proof}

\begin{proposition}\label{ideal_congruence_iff_neg_add}
    Suppose $I$ is a two-sided ideal of $R$.
    Let $x,y\in\carrier[R]$.
    Then $x\mathrel{\idealCongruence{R}{I}}y$ iff
        $\ringAdd[R](\ringNeg[R](x),y)\in I$.
\end{proposition}
\begin{proof}
    Let $G=\AdditiveGroup{R}$.
    $\subgroup{I}{G}$ by
        \obv{two_sided_ideal,left_ideal_additive_subgroup}.
    $x,y\in\carrier[G]$ and
        $\groupInv[G]=\ringNeg[R]$ and
        $\mul[G]=\ringAdd[R]$ by
        \auto{additive_group_fields}.
    Follows by
        \auto{ideal_congruence,subgroup_coset_relation_iff}.
\end{proof}

\begin{proposition}\label{ideal_congruence_iff_add_neg}
    Suppose $I$ is a two-sided ideal of $R$.
    Let $x,y\in\carrier[R]$.
    Then $x\mathrel{\idealCongruence{R}{I}}y$ iff
        $\ringAdd[R](x,\ringNeg[R](y))\in I$.
\end{proposition}
\begin{proof}
    Let $E=\idealCongruence{R}{I}$.
    $\idealCongruence{R}{I}$ is an equivalence on $\carrier[R]$ by
        \obv{ideal_congruence_equivalence}.
    Hence $E$ is an equivalence on $\carrier[R]$.
    $\ringNeg[R](x),\ringNeg[R](y)\in\carrier[R]$ by
        \obv{ring_neg_type}.
    Show that if $x\mathrel{E}y$, then
        $\ringAdd[R](x,\ringNeg[R](y))\in I$.
    \begin{subproof}
        Assume $x\mathrel{E}y$.
        Hence $y\mathrel{E}x$ by \auto{symmetric}.
        Hence $\ringAdd[R](\ringNeg[R](y),x)\in I$ by
            \obv{ideal_congruence_iff_neg_add}.
        $\ringAdd[R](\ringNeg[R](y),x)=
            \ringAdd[R](x,\ringNeg[R](y))$ by
            \obv{ring_add_commutative}.
    \end{subproof}
    Show that if $\ringAdd[R](x,\ringNeg[R](y))\in I$, then
        $x\mathrel{E}y$.
    \begin{subproof}
        Assume $\ringAdd[R](x,\ringNeg[R](y))\in I$.
        $\ringAdd[R](x,\ringNeg[R](y))=
            \ringAdd[R](\ringNeg[R](y),x)$ by
            \obv{ring_add_commutative}.
        Hence $\ringAdd[R](\ringNeg[R](y),x)\in I$.
        Hence $y\mathrel{E}x$ by
            \obv{ideal_congruence_iff_neg_add}.
        Follows by \auto{symmetric}.
    \end{subproof}
    Follows by \auto.
\end{proof}

\begin{proposition}\label{ideal_congruence_reflexive}
    Suppose $I$ is a two-sided ideal of $R$.
    Let $x\in\carrier[R]$.
    Then $x\mathrel{\idealCongruence{R}{I}}x$.
\end{proposition}
\begin{proof}
    $\idealCongruence{R}{I}$ is an equivalence on $\carrier[R]$ by
        \obv{ideal_congruence_equivalence}.
    Follows by \obv{reflexive_on}.
\end{proof}

\begin{proposition}\label{ideal_congruence_symmetric}
    Suppose $I$ is a two-sided ideal of $R$.
    Let $x,y\in\carrier[R]$.
    Suppose $x\mathrel{\idealCongruence{R}{I}}y$.
    Then $y\mathrel{\idealCongruence{R}{I}}x$.
\end{proposition}
\begin{proof}
    $\idealCongruence{R}{I}$ is an equivalence on $\carrier[R]$ by
        \obv{ideal_congruence_equivalence}.
    Follows by \auto{symmetric}.
\end{proof}

\begin{proposition}\label{ideal_congruence_transitive}
    Suppose $I$ is a two-sided ideal of $R$.
    Let $x,y,z\in\carrier[R]$.
    Suppose $x\mathrel{\idealCongruence{R}{I}}y$ and
        $y\mathrel{\idealCongruence{R}{I}}z$.
    Then $x\mathrel{\idealCongruence{R}{I}}z$.
\end{proposition}
\begin{proof}
    $\idealCongruence{R}{I}$ is an equivalence on $\carrier[R]$ by
        \obv{ideal_congruence_equivalence}.
    Follows by \obv{transitive}.
\end{proof}

\begin{theorem}\label{ideal_congruence_addition_congruent}
    Suppose $I$ is a two-sided ideal of $R$.
    Let $x,x',y,y'\in\carrier[R]$.
    Suppose $x\mathrel{\idealCongruence{R}{I}}x'$ and
        $y\mathrel{\idealCongruence{R}{I}}y'$.
    Then
        $\ringAdd[R](x,y)
            \mathrel{\idealCongruence{R}{I}}
        \ringAdd[R](x',y')$.
\end{theorem}
\begin{proof}
    Let $G=\AdditiveGroup{R}$.
    $G$ is a commutative group by
        \obv{two_sided_ideal_ring,
            additive_group_is_exact_commutative_group}.
    $\subgroup{I}{G}$ by
        \obv{two_sided_ideal,left_ideal_additive_subgroup}.
    Hence $\normalSubgroup{I}{G}$ by
        \obv{commutative_group_subgroup_is_normal}.
    $\carrier[G]=\carrier[R]$ and
        $\mul[G]=\ringAdd[R]$ by \auto{additive_group_fields}.
    $\idealCongruence{R}{I}=\subgroupCosetRelation{G}{I}$ by
        \obv{ideal_congruence}.
    Follows by \auto{normal_coset_relation_mul_congruent}.
\end{proof}

\begin{theorem}\label{ideal_congruence_negation_congruent}
    Suppose $I$ is a two-sided ideal of $R$.
    Let $x,y\in\carrier[R]$.
    Suppose $x\mathrel{\idealCongruence{R}{I}}y$.
    Then
        $\ringNeg[R](x)
            \mathrel{\idealCongruence{R}{I}}
        \ringNeg[R](y)$.
\end{theorem}
\begin{proof}
    Let $G=\AdditiveGroup{R}$.
    $G$ is a commutative group by
        \obv{two_sided_ideal_ring,
            additive_group_is_exact_commutative_group}.
    $\subgroup{I}{G}$ by
        \obv{two_sided_ideal,left_ideal_additive_subgroup}.
    Hence $\normalSubgroup{I}{G}$ by
        \obv{commutative_group_subgroup_is_normal}.
    $\carrier[G]=\carrier[R]$ and
        $\groupInv[G]=\ringNeg[R]$ by \auto{additive_group_fields}.
    $\idealCongruence{R}{I}=\subgroupCosetRelation{G}{I}$ by
        \obv{ideal_congruence}.
    Follows by \auto{normal_coset_relation_inverse_congruent}.
\end{proof}

\begin{lemma}\label{ideal_coset_relation_left_multiplication}
    Suppose $I$ is a two-sided ideal of $R$.
    Let $a,b,b'\in\carrier[R]$.
    Let $G=\AdditiveGroup{R}$ and $E=\subgroupCosetRelation{G}{I}$.
    Suppose $b\mathrel{E}b'$.
    Then $\ringMul[R](a,b)\mathrel{E}\ringMul[R](a,b')$.
\end{lemma}
\begin{proof}
    $R$ is a ring by \obv{two_sided_ideal_ring}.
    Let $d=\ringAdd[R](\ringNeg[R](b),b')$.
    $d\in I$ by
        \auto{subgroup_coset_relation_iff,additive_group_fields}.
    Hence $d\in\carrier[R]$ by
        \auto{two_sided_ideal_subseteq_carrier,subseteq}.
    $\ringMul[R](a,d)\in I$ by
        \auto{two_sided_ideal,left_ideal_expanded_iff}.
    $\ringNeg[R](b)\in\carrier[R]$ by \obv{ring_neg_type}.
    $\ringMul[R](a,d)=
        \ringMul[R](a,\ringAdd[R](\ringNeg[R](b),b'))$.
    $\ringMul[R](a,\ringAdd[R](\ringNeg[R](b),b'))=
        \ringAdd[R]
            (\ringMul[R](a,\ringNeg[R](b)),
             \ringMul[R](a,b'))$ by
        \obv{ring_distributive_left}.
    $\ringMul[R](a,\ringNeg[R](b))
        =\ringNeg[R](\ringMul[R](a,b))$ by
        \obv{ring_neg_product_right}.
    Hence $\ringMul[R](a,d)=
        \ringAdd[R]
            (\ringNeg[R](\ringMul[R](a,b)),
             \ringMul[R](a,b'))$.
    Hence
        $\ringAdd[R]
            (\ringNeg[R](\ringMul[R](a,b)),
             \ringMul[R](a,b'))\in I$.
    Follows by
        \auto{subgroup_coset_relation_iff,additive_group_fields,
            ring_mul_type}.
\end{proof}

\begin{lemma}\label{ideal_coset_relation_right_multiplication}
    Suppose $I$ is a two-sided ideal of $R$.
    Let $a,a',b\in\carrier[R]$.
    Let $G=\AdditiveGroup{R}$ and $E=\subgroupCosetRelation{G}{I}$.
    Suppose $a\mathrel{E}a'$.
    Then $\ringMul[R](a,b)\mathrel{E}\ringMul[R](a',b)$.
\end{lemma}
\begin{proof}
    $R$ is a ring by \obv{two_sided_ideal_ring}.
    Let $d=\ringAdd[R](\ringNeg[R](a),a')$.
    $d\in I$ by
        \auto{subgroup_coset_relation_iff,additive_group_fields}.
    Hence $d\in\carrier[R]$ by
        \auto{two_sided_ideal_subseteq_carrier,subseteq}.
    $\ringMul[R](d,b)\in I$ by
        \auto{two_sided_ideal,right_ideal_expanded_iff}.
    $\ringNeg[R](a)\in\carrier[R]$ by \obv{ring_neg_type}.
    $\ringMul[R](d,b)=
        \ringMul[R](\ringAdd[R](\ringNeg[R](a),a'),b)$.
    $\ringMul[R](\ringAdd[R](\ringNeg[R](a),a'),b)=
        \ringAdd[R]
            (\ringMul[R](\ringNeg[R](a),b),
             \ringMul[R](a',b))$ by
        \obv{ring_distributive_right}.
    $\ringMul[R](\ringNeg[R](a),b)
        =\ringNeg[R](\ringMul[R](a,b))$ by
        \obv{ring_neg_product_left}.
    Hence $\ringMul[R](d,b)=
        \ringAdd[R]
            (\ringNeg[R](\ringMul[R](a,b)),
             \ringMul[R](a',b))$.
    Hence
        $\ringAdd[R]
            (\ringNeg[R](\ringMul[R](a,b)),
             \ringMul[R](a',b))\in I$.
    Follows by
        \auto{subgroup_coset_relation_iff,additive_group_fields,
            ring_mul_type}.
\end{proof}

\begin{theorem}\label{ideal_coset_relation_multiplication_congruent}
    Suppose $I$ is a two-sided ideal of $R$.
    Let $a,a',b,b'\in\carrier[R]$.
    Let $G=\AdditiveGroup{R}$ and $E=\subgroupCosetRelation{G}{I}$.
    Suppose $a\mathrel{E}a'$ and $b\mathrel{E}b'$.
    Then $\ringMul[R](a,b)\mathrel{E}\ringMul[R](a',b')$.
\end{theorem}
\begin{proof}
    $I$ is a left ideal of $R$ by \obv{two_sided_ideal}.
    Hence $\subgroup{I}{G}$ by \obv{left_ideal_additive_subgroup}.
    Hence $E$ is an equivalence on $\carrier[G]$ by
        \auto{subgroup_coset_relation_equivalence}.
    $\ringMul[R](a,b)\mathrel{E}\ringMul[R](a',b)$ by
        \obv{ideal_coset_relation_right_multiplication}.
    $\ringMul[R](a',b)\mathrel{E}\ringMul[R](a',b')$ by
        \obv{ideal_coset_relation_left_multiplication}.
    Follows by \obv{transitive}.
\end{proof}

\begin{theorem}\label{ideal_congruence_quotient_projection_eq_iff}
    Suppose $I$ is a two-sided ideal of $R$.
    Let $x,y\in\carrier[R]$.
    Let $G=\AdditiveGroup{R}$ and
        $q=\groupQuotientProjection{G}{I}$.
    Then $q(x)=q(y)$ iff $x\mathrel{\idealCongruence{R}{I}}y$.
\end{theorem}
\begin{proof}
    Let $E=\idealCongruence{R}{I}$.
    $G$ is a commutative group by
        \obv{two_sided_ideal_ring,
            additive_group_is_exact_commutative_group}.
    $\subgroup{I}{G}$ by
        \obv{two_sided_ideal,left_ideal_additive_subgroup}.
    Hence $\normalSubgroup{I}{G}$ by
        \obv{commutative_group_subgroup_is_normal}.
    $\idealCongruence{R}{I}$ is an equivalence on $\carrier[R]$ by
        \obv{ideal_congruence_equivalence}.
    Hence $E$ is an equivalence on $\carrier[R]$.
    $\carrier[G]=\carrier[R]$ and
        $E=\subgroupCosetRelation{G}{I}$ by
        \auto{additive_group_fields,ideal_congruence}.
    $q(x)=\equivalenceClass{E}{x}$ and
        $q(y)=\equivalenceClass{E}{y}$ by
        \auto{group_quotient_projection_apply}.
    Follows by \auto{equiv_iff_equivclasses_eq}.
\end{proof}
